\chapter{Conclusiones y Recomendaciones}

\section{Conclusiones}

\begin{itemize}
    \item \textbf{Sobre la implementación algorítmica:} se implementaron desde cero, sin librerías de optimización ni de grafos, los cuatro paradigmas exigidos y aplicados a los casos de uso del enunciado: Greedy para la ruta voraz y la asignación de emergencias; Divide y Vencerás para segmentar la ciudad en una zona por repartidor; Programación Dinámica para cargar el vehículo (Mochila 0/1) y obtener el recorrido exacto (Held-Karp); y Backtracking para encontrar todas las rutas de San Blas a Wanchaq evitando calles cerradas. Su correctitud está respaldada por 39 pruebas unitarias con resultados esperados calculados de forma independiente.

    \item \textbf{Sobre la complejidad:} las mediciones concuerdan con el análisis Big-O. El tiempo de Held-Karp, $O(n^2 2^n)$, se multiplicó por 376 entre 5 y 12 puntos, mientras que el Greedy, $O(n^2)$, se mantuvo por debajo de 0.05 ms. El análisis también corrigió dos estimaciones iniciales: la partición de Divide y Vencerás cuesta $\Theta(n\log^2 n)$, no $\Theta(n\log n)$, porque en cada nivel se ordena; y su variante con un solo repartidor cuesta $O(n^2/K)$ por la elección voraz entre zonas.

    \item \textbf{Sobre la calidad de las soluciones:} el Greedy produjo recorridos hasta 24.5\% más largos que el óptimo, y la mochila voraz cargó hasta S/.~15 menos que la PD. La poda del Backtracking redujo los estados visitados entre 42\% y 86\% sin perder la ruta óptima. Dividir la ciudad en una zona por repartidor redujo la distancia total de 43.4 a 30.2 km.

    \item \textbf{Sobre el sistema:} la aplicación ejecuta todas sus funcionalidades sin bloquearse, gracias a los cálculos en segundo plano y al manejo centralizado de errores; funciona sin conexión a Internet con un mapa vectorial propio; y permite cerrar calles y observar su efecto en todos los algoritmos. Una prueba automática de la interfaz recorre todas las pestañas y respalda esta afirmación.
\end{itemize}

\section{Recomendaciones}

\begin{itemize}
    \item Para escalar a más de 14 paradas, sustituir Held-Karp por heurísticas de mejora como 2-opt aplicadas sobre la ruta del Greedy, que reducen el exceso de distancia en tiempo polinomial.
    \item Incorporar el volumen como segunda restricción de la mochila, con una tabla $dp[i][w][v]$ de costo $O(N\cdot W\cdot V)$; actualmente el volumen solo se verifica después de optimizar por peso.
    \item Calcular la mediana por selección lineal (Quickselect) en lugar de ordenar en cada nivel, para que la partición espacial baje de $\Theta(n\log^2 n)$ a $\Theta(n\log n)$.
    \item Revisar en campo los ocho lugares cuya ubicación se desplazó más de 60 m al ajustarse a la red vial (Anexo B) y reconstruir el caché vial con las coordenadas corregidas.
    \item Añadir persistencia de pedidos (por ejemplo, SQLite) y tiempos de viaje por tipo de vehículo, para que la herramienta pueda usarse en la operación diaria de una empresa de reparto.
\end{itemize}
